%%%%%%%%%%%%%%%%%%%%%%%%%%%%%%%%%%%%%%%%%%%%%%%%%%%%%%%%%%%%%%%%%%
%% Toto je inkludovaný LaTeXový soubor `c_intro.tex'.           %%
%% ------------------------------------------------------------ %%
%% Verze 2022-05-31                                             %%
%% Vyvíjí & udržuje <stanislav.hledik@physics.slu.cz>           %%
%%%%%%%%%%%%%%%%%%%%%%%%%%%%%%%%%%%%%%%%%%%%%%%%%%%%%%%%%%%%%%%%%%

%% Text vaší závěrečné práce by měl začínat nečíslovanou kapitolou
%% Úvod. Potlačení číslování lze dosáhnout pomocí "hvězdičkové
%% formy" příkazu \chapter, jak je vidět níže. Dokumentní třída
%% je naprogramována tak, aby do Obsahu vložila i nečíslované
%% kapitoly. Všimněte si, jak je každý příkaz, který generuje
%% nějaký čítač, označen unikátním štítkem pomocí \label. To se
%% může hodit pro pozdější křížové reference.

%#################################################################
\chapter*{Úvod}\label{intro}

%% Motto, pokud je chcete:

\begin{motto}{Murphy (Allenův nebo Cannův axiom)}
Když už nefunguje vůbec nic, přečti si konečně návod k~použití.
\end{motto}

%% Implicitně je první odstavec po nadpisu (kapitola, sekce, ...)
%% neindentovaný (bez horizontální odstavcové zarážky, to lze ale
%% globálně změnit aktivací volby `indentfirst'). V tomto případě
%% se však nejedná o první odstavec, jelikož před ním se vyskytuje
%% prostředí `motto'. Proto musíme indentaci explicitně potlačit.
%% Tento odstavec rovněž demostruje vložení hesla pro rejstřík:
\noindent
Publikační systém\INDEX{publikační systém} \LaTeX{}\INDEX{publikační
  systém!\LaTeX{}}~\cite{Lam:1994:LaTeX:,%
  Goo-Mit-Sam:1994:Companion:,%
  LaTeXNet:LIVE:TheLaTeXNetwork:} je de~facto standardem pro sazbu prací
obsahujících množství matematiky. Do této kategorie dokumentů spadají také
bakalářské, diplomové, rigorózní a dizertační práce (jež jsou v~dalším
souhrnně nazývány \emph{závěrečnými pracemi}) v~matematických, fyzikálních a
informatických oborech.

%% Obsahuje-li Úvod sekce hlubší úrovně (\section, ...), potlačte
%% jejich číslování "hvězdičkovou formou":
\section*{Cíl}\label{aim}

Po ustavení Fyzikálního ústavu v~Opavě v~r.~2020 vyvstala v~návaznosti na nové
studijní předpisy a množící se dotazy studentů na formální uspořádání
závěrečné práce potřeba poskytnout vhodné \LaTeX ové řešení.  Po rešerši
balíčků dostupných z~CTAN~\cite{CTAN:LIVE:CompTeXArNet:} a
z~Overleaf~\cite{Ovr:LIVE::Overleaf} se autor rozhodl nové řešení založit na
předchozích balíčcích pro Ústav fyziky FPF SU v~Opavě \pkg{physuth} (z~let
2007--2009) a \pkg{fpfslut} (2015), avšak pod změněným názvem \pkg{\ipotname}
(akronym Institute of Physics in Opava THESis).

Klíčovými požadavky při vývoji balíčku \pkg{\ipotname} byla
\begin{enumerate}
\item\label{easyuse} snadnost instalace a základního použití pro ty, kdo
  nemají zkušenosti s~publikačním systémem \LaTeX;
\item\label{autotitle} snadné generování obsahově a graficky jednotných
  titulních stránek\footnote{Sestávajících z~patitulu, titulní stránky,
    abstraktu, klíčových slov, zadání práce, prohlášení, obsahu, volitelně
    protititulu (frontispice), poděkování, dedikace, seznamu obrázků a seznamu
    tabulek.} (\uv{front matter});
\item\label{langsup} podpora jazyků schválených pro psaní závěrečných prací na
  Fyzikálním ústavu;
\item\label{bibgen} podpora tvorby bibliografie ve formě bibliografického
  stylu\INDEX{bibliografický styl} pro pomocný program
  \BibTeX{}~\cite{Pat:LIVE:CTAN:BibTeX};
\item\label{indgen} podpora tvorby rejstříku pomocí pomocného programu
  \MakeIndex{}~\cite{Bee:LIVE:CTAN:MakeIndex,%
    LtX:LIVE:CTAN:makeidx};
\item\label{advposs} konfigurovatelnost podle potřeb pokročilejšího uživatele;
\item\label{sufdoc} dostatečná a srozumitelná dokumentace.
\end{enumerate}

\section*{Prerekvizity}\label{prereq}

Předpokládá se, že čtenář je obeznámen se základy publikačního
systému\INDEX{publikační systém} \LaTeX\INDEX{publikační systém!\LaTeX{}}
zhruba v~rozsahu volitelného předmětu bakalářského studijního programu Fyzika
FU:FYBPV0006 \LaTeX{}\footnote{Viz
  \url{https://is.slu.cz/predmet/fu/FYBPV0006}.} nebo
příručky~\cite{Ryb:2005:LaTeX:3}, proto se zabývá jen specifiky \LaTeX
ového\INDEX{publikační systém!\LaTeX{}} balíčku \pkg{\ipotname} a
nesubstituuje běžné manuály~\cite{Ryb:2005:LaTeX:3,%
  Lam:1994:LaTeX:,%
  Goo-Mit-Sam:1994:Companion:,%
  Goo-Rah-Mit:1994:GraComp:,%
  Goo-Rah:1994:WebComp:,%
  Kop-Dal:2004:LaTeXPodrPruv:,%
  Wri-Fai-etal:LIVE::TeXFAQ} nebo online dokumentaci a
příručky~\cite{Oet-etal:LIVE::lshort,%
  Hle:LIVE::Bazak}.  Pro podrobnější poučení o~problematice sazby závěrečných
prací lze doporučit~\cite{Pol:1998:PravidlaSazbyDP:,%
  Ryb:2006:ZpracZP:} a odkazy v~nich uvedené.

\section*{Struktura}\label{struct}

Tento dokument má formu fiktivní bakalářské práce sloužící jako uživatelský
manuál a je rozdělen do tří částí. V~Části~\ref{procsol} jsou soustředěny
popisy konkrétních činností, postupů a řešení (zhruba v~pořadí, v~němž se
s~nimi uživatel setkává), dále tipy a triky uplatnitelné při psaní závěrečných
prací pomocí \pkg{\ipotname}.

Část~\ref{findef} sestává z~poněkud nesourodých ukázek vybraných a mírně
upravených a/nebo zkrácených kapitol a partií ze skutečných již obhájených
závěrečných prací (s~laskavým svolením jejich autorů) s~vesměs zaslepenými
obrázky (to hlavně kvůli redukci velikosti obrázkových souborů). Konkrétně se
jedná o~ukázky z~bakalářských prací~\cite{Val:2015:BcProj:NumMagLev,%
  Bry:2007:BcProj:ModMedDiag} na stránkách~\pageref{SpinLev}
a~\pageref{princtwo}, o~ukázku z~diplomové
práce~\cite{Fab:2017:DipProj:PorSnimkuOpAnObr} na stránce~\pageref{CAPczm} a
o~kapitolu z~vlastní dizertační práce autora balíčku
\pkg{\ipotname}~\cite{Hle:2014:PhDThes:Polytropes} na
stránce~\pageref{Chap:Uniform}.\footnote{Zatímco první tři ukázky jsou psány
  česky, ukázka z~dizertační práce zůstává v~angličtině~-- pro účel, jemuž
  mají ukázky sloužit, je jejich jazyk irelevantní.} Do této části je~-- bez
nároku na úplnost~-- zahrnuta většina typických situací, se kterými se lze při
sazbě závěrečných prací setkat. Čtenářův předpokládaný modus operandi pak
spočívá ve vybrání požadovaného typografického efektu ve vysazené fromě práce
% (třeba obrázku jakožto plovoucího objektu s~určitým aranžmá grafických souborů
% a společnou popiskou, mimo text vysazených očíslovaných rovnic zarovnaných
% podle rovnítka, tabulky s~popiskou)
a následného vyhledání odpovídajícího místa v~hlavním zdrojovém textu
\pgm{\thesnamecze .tex} nebo v~inkludovaných zdrojových textech
\pgm{include/*.tex} a dalších pomocných textových souborech jako jsou
bibliografická databáze \pgm{\refername .bib} či konfigurační soubory
\pgm{\ipotname .cfg} a \pgm{hyperref.cfg}, kde kromě vlastního řešení tu a tam
nalezne další informace o~žádaném efektu formou komentářů.\footnote{Hledáte-li
  např. inspiraci pro aranžmá obrázků nebo tabulek jakožto plovoucích objektů
  a jejich umisťování, lze doporučit ukázku z~bakalářské
  práce~\cite{Bry:2007:BcProj:ModMedDiag}.} Ostatně lze vyjít přímo ze
zdrojových textů této fiktivní bakalářské práce a vlastní práci začít psát
jejich přepisováním textem novým. Tento způsob \uv{učení se podle příkladu}
jistě nepokryje všechny požadavky, ale spolu s~uvedenou literaturou může
poskytnout vodítka a programovací vzory pro vlastní kreativní tvorbu.

Technické detaily (konfigurace, volby dokumentní třídy \pgm{\ipotname .cls},
načítané balíčky, pokročilé způsoby kompilace pomocí \pkg{make} a jiné) jsou
popsány v~dodatkové Části~\ref{append}.

\section*{Další vývoj}\label{furtherdev}%%%%%%%%%%%%%%%%%%%%%%%%%%%%%%%%%%%

Knihovny balíčku \pkg{\ipotname} a hlavně tento manuál budou v~budoucnu
podléhat dalšímu vývoji. Ve svém současném stavu je manuál zatím
minimalistický a bude doplňován o~další plánované i~neplánované kapitoly
(konstrukce preambule, sazba matematiky, tvorba titulních stránek, vkládání
externí grafiky, citace a reference, jazyková lokalizace) a dodatky
(pokročilejší konfigurace, manipulace se sazebním obrazcem, hypertext,
management kompilace pomocí \pkg{make}). Sledujte proto domovskou stránku
odkázanou v~Kapitole~\ref{install}, na níž bude každá nová verze anoncována.

Protože ani balíček \pkg{\ipotname}~-- stejně jako každý jiný software~-- není
prost chyb, chtěl bych požádat uživatele, aby mě pomocí kontaktů uvedených na
\authorurl{} laskavě informovali o~chybách, nedostatcích a problémech, na
které při jeho používání narazí, za což předem děkuji.

\endinput

%% 
%% End of file `c_intro.tex'.
